\documentclass[12pt]{article}
\usepackage{amsmath}
\usepackage{tikz}
\usepackage{soul}
\usepackage[empty]{fullpage}

\begin{document}

\section*{Project 2 - Tips}

\begin{enumerate}
\item To write the functions \verb|letter2number| and \verb|number2letter|, it would help to use the following string. 
\begin{center}
\verb|alphabet = "abcdefghijklmnopqrstuvwxyz"|
\end{center}
Recall that you can use the string method \verb|.find()| to find the index of any letter in the \verb|alphabet| string. 

\item You will definitely need to use modular arithmetic at some point during this project. It might help if your \verb|number2letter| function uses modular arithmetic to handle numbers bigger than 25.  For example, it might be nice if 26 wraps around to 0 and returns the letter \verb|"a"|.  

\item You don't need to worry about capital letters in this project, just write your code to handle lower case letters. 

\item You should only apply the Caesar shift to letter characters, not digits or punctation symbols.  To check if a character is a letter of the alphabet, you could use the \verb|in| operator with the \verb|alphabet| variable, or the \verb|.isalpha()| method.  


%There are at least two ways to do this.  One option is to create a string with all of the letters of the alphabet in order and use the \verb|.index()| method.  Another option is to use the built in functions \verb|ord()| and \verb|chr()|.  The function \verb|ord()| converts any character to its integer ASCII value.  So \verb|ord("a")| returns 97, \verb|ord("b")| returns 98, and so on. The function \verb|chr()| does the reverse, it converts an integer to its corresponding ASCII character.  So \verb|chr(97)| is \verb|"a"|, \verb|chr(98)| is \verb|"b"|, etc. 

%\item Don't forget methods like \verb|.count()| when working with lists and strings.  

%\item For the \verb|argmax()| function, you'll need to loop over every value in the input list.  Make sure to use an accumulator variable to keep track of which index corresponds to the maximum value you have seen so far. 
\end{enumerate}



\end{document}
